\section{Introduction}
\label{sec:intro}

Sepsis is one of the leading causes of death in intensive care units worldwide, accounting for millions of fatalities each year. Treatment is high-stakes and sequential: the clinician decides how much vasopressor to administer to stabilise blood pressure and how much intravenous fluid to give to maintain circulation, then observes the patient's response and decides again. Too little intervention and the patient deteriorates; too much and the treatment itself causes harm. The choice is repeated every few hours over a stay that may last days.

Reinforcement learning is a natural framing for this problem. The clinician's decision rule is a policy, the patient's evolving physiology is a state, and survival is a delayed terminal reward. An RL agent trained on historical trajectories could in principle uncover treatment strategies that improve outcomes beyond what clinical guidelines alone produce. The same framing also makes the limitations visible: the dataset is observational, the action space is coarse, and any policy learned in a simulator may not transfer to bedside care without further validation.

\subsection{The benchmark}

ICU-Sepsis-v2~\citep{choudhary2024icusepsis} is a gymnasium environment fitted to the MIMIC-III critical care database~\citep{johnson2016mimic} using the AI Clinician methodology of \citet{komorowski2018ai}. The MDP has 716 states (714 clinical, 2 terminal corresponding to survival and death) and 25 actions, decomposed as 5 vasopressor levels by 5 IV fluid levels. The reward is sparse: $+1$ on survival, $0$ on death, $0$ at all intermediate steps. Episodes terminate when the patient outcome is realised, with discount factor $\gamma = 1$.

This project applies two configurations of the same MDP. \emph{Configuration A} exposes the discrete state index directly. Tabular RL is feasible at that scale, and the agent has access to the full transition tensor $P$ and reward matrix $R$. \emph{Configuration B} replaces the state index with the 47-dimensional continuous physiological feature vector that produced the state clustering~\citep{komorowski2018ai}, and adds three orthogonal clinical-reality wrappers: episodic observation noise (monitor malfunction), episodic missing labs (analyser downtime), and rare acute events (sudden unrecoverable mortality). The wrappers introduce partial observability and irreducible stochasticity.

A fixed difficulty configuration applies throughout. The initial-state distribution is biased toward sicker patients through $\text{SOFA\_BIAS} = 5.0$, which raises the mean starting SOFA score from approximately 7 to 9. An intensity penalty $\lambda = 0.02$ is subtracted from each step's reward in proportion to the action's combined vasopressor and fluid level. These settings define the cohort and the reward structure against which every algorithm is evaluated.

\subsection{Approach and structure}

The project answers three questions, in order. First, how close does tabular Q-Learning come to the model-based optimum on the discrete MDP, and what does the residual gap reveal about the learning algorithm? Second, what does the same algorithmic family achieve when the observation becomes a noisy continuous vector that requires function approximation? Third, can a risk-sensitive reformulation of the objective produce a clinically more conservative policy without losing aggregate performance?

\Cref{sec:methodology} describes the algorithms and their hyperparameters. \Cref{sec:evaluation} presents the comparative analysis across both configurations, including a bridge experiment that decomposes the function-approximation cost from the wrapper cost. \Cref{sec:extension} introduces the risk-sensitive extension. \Cref{sec:conclusion} summarises the findings and lists limitations.

All numerical results in this report come from a single notebook deliverable that produces the figures and tables reproducibly under fixed seeds. The accompanying source code is in the RF\_group\_X repository; the notebook is \texttt{rl\_sepsis\_project\_refactored.ipynb}.
